\chapter{El caso PETROPERÚ: una empresa expuesta}\label{cap:caso}

\begin{objetivos}
Al terminar este capítulo usted podrá:
\begin{itemize}
\item Describir el negocio de PETROPERÚ y de dónde vienen sus ingresos y sus costos.
\item Leer los indicadores del Cuadro~1 del informe: ingresos, resultado neto, capital de trabajo, pasivos financieros y apalancamiento.
\item Explicar qué es la moneda funcional y por qué genera riesgo cambiario.
\item Interpretar una calificación crediticia y la expresión \enquote{incertidumbre material sobre empresa en marcha}.
\item Clasificar los riesgos financieros en riesgo de mercado, de crédito y de liquidez.
\item Explicar por qué una empresa con poca caja tiene más razones para cubrirse.
\end{itemize}
\end{objetivos}

\begin{prerreq}Ninguno. Este capítulo es la puerta de entrada.\end{prerreq}

\section{Por qué empezar por una empresa}

Imagine que usted administra una bodega de barrio. Compra arroz al mayorista y lo vende a sus vecinos. Si un mes el precio mayorista sube mucho y usted no puede subir de inmediato el precio de venta, gana menos o pierde. Si usted además debe un préstamo en una moneda distinta a la de sus ventas, su deuda puede crecer aunque usted no haya pedido un sol más. Dos variables que usted no controla, el precio de lo que compra y el tipo de cambio, deciden si a fin de mes le alcanza para pagar.

PETROPERÚ enfrenta exactamente este problema, a una escala de miles de millones de dólares. Los derivados financieros son, en esencia, contratos que permiten \emph{trasladar} parte de esa incertidumbre a alguien dispuesto a asumirla. Antes de estudiar los contratos, necesitamos entender bien la incertidumbre que se quiere trasladar.

\section{El negocio de PETROPERÚ}

PETROPERÚ es la empresa estatal que refina y comercializa combustibles en el Perú. Opera la Nueva Refinería Talara (NRT), las refinerías de Iquitos y Conchán, y el Oleoducto Norperuano. Su cadena de valor puede resumirse en tres pasos:

\begin{enumerate}
\item \textbf{Compra} petróleo crudo y productos. En 2025 compró por US\$ 2,490.0 millones, a un precio promedio de US\$ 78.66 por barril y con un volumen de 87 MBDC (miles de barriles por día calendario). Los precios del crudo se fijan en dólares y siguen referencias internacionales como el WTI o el Brent.
\item \textbf{Refina} ese crudo y lo transforma en gasolinas, diésel, petróleo industrial y otros productos. Entre la compra del crudo y la venta del producto refinado pasan varias semanas.
\item \textbf{Vende} los combustibles en el mercado local. Se facturan en soles, pero sus precios siguen la \emph{paridad de importación}: lo que costaría importar ese mismo combustible, que se fija en dólares.
\end{enumerate}

\begin{definicion}[Activo subyacente y precio de referencia]
Un \textbf{precio de referencia} (o \emph{benchmark}) es un precio público y ampliamente observado que sirve para fijar el precio de muchas transacciones. El \textbf{WTI} (\emph{West Texas Intermediate}) es un tipo de crudo cuyo precio se negocia en la bolsa NYMEX de Nueva York; el \textbf{Brent} es la referencia del Mar del Norte. Cuando un contrato financiero depende de un precio así, ese precio (o el bien que lo origina) se llama \textbf{activo subyacente} del contrato.
\end{definicion}

\section{Leer los estados financieros}

El informe se apoya en los estados financieros auditados al 31 de diciembre de 2025. Antes de interpretar sus cifras, conviene repasar qué mide cada estado.

\begin{definicion}[Estados financieros básicos]
\begin{itemize}
\item El \textbf{estado de situación financiera} (o balance) es una foto en una fecha: lo que la empresa tiene (\emph{activos}), lo que debe (\emph{pasivos}) y la diferencia (\emph{patrimonio}). Siempre se cumple $\text{Activo}=\text{Pasivo}+\text{Patrimonio}$.
\item El \textbf{estado de resultados} es una película de un período: ingresos menos costos y gastos da el resultado (utilidad o pérdida).
\item Los activos y pasivos \textbf{corrientes} son los que se convierten en caja o se pagan dentro de los próximos doce meses.
\end{itemize}
\end{definicion}

El Cuadro~1 del informe resume los indicadores clave. Lo reproducimos con una explicación de cada línea.

\begin{center}\small
\captionof{table}{Indicadores financieros de PETROPERÚ (US\$ millones). Fuente: Cuadro~1 del informe.}\label{tab:c1}
\begin{tabularx}{\linewidth}{@{}>{\raggedright\arraybackslash}Xrr@{}}
\toprule
& 2024 & 2025 \\
\midrule
Ingresos de actividades ordinarias & 3,456.4 & 3,380.4\\
Resultado neto & (773.9) & (601.5)\\
Capital de trabajo & (1,900.0) & (2,296.4)\\
Otros pasivos financieros (100\,\% a tasa fija) & 5,610.7 & 5,168.8\\
Deuda neta / capital total & 0.70 & 0.78\\
\bottomrule
\end{tabularx}
\end{center}

\paragraph{Ingresos} Son las ventas del año. Cayeron levemente en 2025 por el menor precio del crudo: como los precios de venta siguen la paridad de importación, cuando el crudo baja también bajan los ingresos.

\paragraph{Resultado neto} Es la última línea del estado de resultados. PETROPERÚ perdió US\$ 601.5 millones en 2025, menos que los US\$ 773.9 millones de 2024, pero sigue perdiendo.

\paragraph{Utilidad bruta} No está en el cuadro pero el informe la menciona: fue de apenas US\$ 14.4 millones. La utilidad bruta es ingresos menos costo de ventas, antes de gastos administrativos, financieros e impuestos. Comparada con ingresos de US\$ 3,380.4 millones, representa un margen bruto de $14.4/3{,}380.4 = 0.43\,\%$. Por cada 100 dólares vendidos, quedan 43 centavos para cubrir todo lo demás.

\paragraph{Capital de trabajo} Se define como activo corriente menos pasivo corriente. Un valor negativo de US\$ 2,296.4 millones significa que, si hubiera que pagar hoy todo lo que vence en un año, los activos que se convierten en caja en ese plazo no alcanzarían por esa diferencia. Es la señal más directa de estrechez de liquidez.

\paragraph{Otros pasivos financieros} Son las deudas con bancos y bonistas. El dato clave del cuadro es la nota \enquote{100\,\% a tasa fija}: la empresa sabe cuánto interés pagará; su riesgo de tasas, como veremos en el capítulo~\ref{cap:riesgo}, está en otro lugar.

\paragraph{Deuda neta / capital total} La deuda neta es la deuda financiera menos la caja. El capital total es deuda neta más patrimonio. Un ratio de 0.78 indica que el 78\,\% del financiamiento de largo plazo de la empresa proviene de acreedores y solo el 22\,\% de los accionistas: la empresa está muy \textbf{apalancada}.

\begin{ejemplo}{Capital de trabajo}
Suponga una empresa con activo corriente de 500 y pasivo corriente de 800 (en millones). Su capital de trabajo es $500-800=-300$. Si en los próximos doce meses debe pagar 800 y solo cobrará o liquidará 500, le faltan 300 que deberá conseguir con nuevos préstamos, aportes de capital o ventas de activos de largo plazo. PETROPERÚ está en esta situación con un faltante de US\$ 2,296.4 millones.
\end{ejemplo}

\subsection{La moneda funcional}

\begin{definicion}[Moneda funcional]
La \textbf{moneda funcional} es la moneda del entorno económico principal en el que opera la empresa: aquella en la que se fijan sus precios de venta y sus principales costos. Los estados financieros se miden en esa moneda; todo lo demás es \enquote{moneda extranjera} y debe convertirse.
\end{definicion}

La moneda funcional de PETROPERÚ es el dólar estadounidense, aunque la empresa sea peruana y facture en soles. La razón es económica: sus compras de crudo se pactan en dólares y sus precios de venta siguen la paridad de importación en dólares. Consecuencia importante: \emph{para PETROPERÚ, el sol es una moneda extranjera}. Cualquier deuda en soles se mide en dólares, y su valor en dólares cambia cuando cambia el tipo de cambio.

\begin{intuicion}[deber soles cuando uno piensa en dólares]
Si usted debe S/ 3,400 y el tipo de cambio es 3.40 soles por dólar, su deuda equivale a US\$ 1,000. Si el sol se aprecia y el tipo de cambio baja a 3.10, la misma deuda de S/ 3,400 equivale a $3{,}400/3.10=\text{US\$}~1{,}096.8$. Usted no pidió prestado nada adicional, pero en su moneda de medida debe casi 97 dólares más. Esta es exactamente la pérdida cambiaria que el informe cuantifica.
\end{intuicion}

\subsection{Empresa en marcha y calificación crediticia}

El auditor incluyó en su dictamen una \textbf{incertidumbre material sobre la continuidad como empresa en marcha}. Los estados financieros se preparan bajo el supuesto de que la empresa seguirá operando; cuando el auditor advierte una incertidumbre material sobre ese supuesto, está diciendo que hay dudas significativas sobre la capacidad de la empresa de seguir funcionando sin apoyo extraordinario (por ejemplo, del Estado).

\begin{definicion}[Calificación crediticia]
Una \textbf{calificación crediticia} es una opinión emitida por una agencia especializada (S\&P, Moody's, Fitch) sobre la probabilidad de que un emisor pague sus deudas a tiempo. Las calificaciones desde BBB- (S\&P) o Baa3 (Moody's) hacia arriba se llaman \textbf{grado de inversión}; por debajo, \textbf{grado especulativo}. La \textbf{perspectiva} (positiva, estable, negativa) indica hacia dónde podría moverse la calificación.
\end{definicion}

\begin{informe}
S\&P rebajó a PETROPERÚ a B- en diciembre de 2025 y la reafirmó en mayo de 2026 con perspectiva negativa; Moody's la reafirmó en Caa1 en setiembre de 2026 con perspectiva positiva; Fitch retiró sus calificaciones en enero de 2026. Las calificaciones B- y Caa1 están varios escalones por debajo del grado de inversión.
\end{informe}

¿Por qué importa la calificación en un libro sobre derivados? Porque un derivado es un contrato con obligaciones futuras para ambas partes. Un banco que firma un derivado con PETROPERÚ asume el riesgo de que la empresa no cumpla. Cuanto peor la calificación, más le cobrará el banco por ese riesgo o más garantías le pedirá. Esta idea, que el informe llama CVA y colateral, se desarrolla en el capítulo~\ref{cap:mercados}.

\section{¿Qué es riesgo financiero?}

En el lenguaje cotidiano, \enquote{riesgo} es la posibilidad de que algo salga mal. En finanzas el término es más preciso.

\begin{definicion}[Riesgo financiero]
\textbf{Riesgo} es la posibilidad de que un resultado futuro (un flujo de caja, una utilidad, el valor de un activo) se desvíe de lo esperado, en cualquier dirección. Una \textbf{exposición} es la cantidad que está sujeta a ese riesgo: el monto cuyo valor cambia cuando cambia una variable de mercado.
\end{definicion}

Observe que la definición incluye desviaciones favorables. Un precio del crudo que sube puede beneficiar el inventario; un precio que baja lo perjudica. El riesgo es la \emph{variabilidad}, no solo la pérdida. Esta simetría será importante: cubrirse suele significar renunciar a las sorpresas buenas para eliminar las malas.

La Nota~3.1 de los estados financieros, y el informe con ella, clasifica los riesgos así:

\begin{itemize}
\item \textbf{Riesgo de mercado}: pérdidas por cambios en precios de mercado. Incluye:
  \begin{itemize}
  \item \emph{riesgo de tipo de cambio} (el precio de una moneda en términos de otra);
  \item \emph{riesgo de tasa de interés} (el precio del dinero en el tiempo);
  \item \emph{riesgo de precio de materias primas} (aquí, el crudo).
  \end{itemize}
\item \textbf{Riesgo de crédito}: pérdidas porque un deudor o contraparte no paga.
\item \textbf{Riesgo de liquidez}: no tener caja suficiente cuando hay que pagar, aunque la empresa tenga activos.
\end{itemize}

\begin{importante}[Riesgo de precio frente a restricción de liquidez]
El informe hace una distinción fina que conviene retener desde ahora: el riesgo de liquidez de PETROPERÚ \emph{no se cubre con derivados}, porque no es una variable de mercado que suba o baje. Pero \emph{condiciona el diseño} de las coberturas: no se pueden usar instrumentos que exijan pagos iniciales grandes (primas) o depósitos diarios de garantía (márgenes).
\end{importante}

\section{¿Por qué una empresa debería cubrirse?}

Esta pregunta tiene una respuesta menos obvia de lo que parece.

\begin{adicional}[el argumento de la irrelevancia]
En un mundo ideal sin impuestos, sin costos de quiebra y donde la empresa puede pedir prestado siempre a una tasa justa, cubrir riesgos no crea valor para los accionistas: ellos mismos podrían diversificarse o cubrirse por su cuenta. Es una aplicación de las ideas de Modigliani y Miller. Por eso la pregunta correcta no es \enquote{¿reduce riesgo?} sino \enquote{¿qué imperfección del mundo real hace que reducir riesgo valga la pena?}.
\end{adicional}

Froot, Scharfstein y Stein (1993) dieron una respuesta influyente: el financiamiento externo es caro, sobre todo para empresas en dificultades. Si una caída de precios vacía la caja justo cuando la empresa necesita invertir (en mantenimiento, en inventario), deberá endeudarse a una tasa muy alta o dejar de invertir. La cobertura sirve para que la caja interna esté disponible cuando se la necesita. Stulz (1996) añade que el objetivo razonable de una empresa no financiera es eliminar las \emph{malas noticias extremas} que la acercarían a la quiebra, no especular sobre la dirección de los precios.

\begin{ejemplo}{Por qué la caja estable vale más (didáctico)}
Una empresa necesita invertir 100 el próximo año en mantenimiento indispensable. Su caja operativa será 130 o 70, con igual probabilidad, según el precio del crudo. El valor esperado de la caja es $0.5\times130+0.5\times70=100$.

\emph{Sin cobertura.} En el escenario bueno invierte 100 y le sobran 30. En el malo le faltan 30, que debe pedir prestados con un recargo de 25\,\% por ser una empresa en dificultades: el costo adicional es $0.25\times30=7.5$. Costo esperado de financiamiento: $0.5\times 0+0.5\times7.5=3.75$.

\emph{Con cobertura.} Un contrato que fija el precio hace que la caja sea 100 en ambos escenarios. Invierte 100 sin pedir prestado. Costo esperado de financiamiento: 0.

La cobertura no cambió la caja esperada (100 en ambos casos), pero ahorró en promedio 3.75 de costos de financiamiento. Ese ahorro es la fuente de valor de la cobertura.
\end{ejemplo}

\begin{informe}
El informe cita exactamente este razonamiento: \enquote{cada dólar que se pierde por tipo de cambio o por el precio del crudo reduce una liquidez escasa} (Froot et al., 1993), y concluye que el objetivo de las coberturas es \enquote{estabilizar resultados y flujos, no especular} (Stulz, 1996).
\end{informe}

\section{Primer contacto con los derivados}

\begin{definicion}[Derivado financiero]
Un \textbf{derivado} es un contrato cuyo valor \emph{deriva} del valor de otra variable, el activo subyacente: un tipo de cambio, una tasa de interés, el precio del petróleo, el de una acción. El contrato establece hoy qué pagos se harán en el futuro en función de lo que ocurra con esa variable.
\end{definicion}

Hay cuatro familias básicas, que estudiaremos en la Parte~II:

\begin{itemize}
\item \textbf{Forward}: acuerdo para comprar o vender algo en una fecha futura a un precio fijado hoy (cap.~\ref{cap:forwards}).
\item \textbf{Futuro}: un forward estandarizado que se negocia en una bolsa y se liquida diariamente (cap.~\ref{cap:futuros}).
\item \textbf{Swap}: intercambio periódico de flujos, por ejemplo intereses en soles por intereses en dólares (cap.~\ref{cap:swaps}).
\item \textbf{Opción}: el \emph{derecho}, pero no la obligación, de comprar o vender a un precio fijado (caps.~\ref{cap:opciones} y~\ref{cap:valoracion}).
\end{itemize}

\begin{informe}[qué hace hoy PETROPERÚ]
La Nota~3 de los estados financieros confirma que la empresa no usa derivados para cubrir su riesgo cambiario. El único derivado revelado son swaps de monedas de corto plazo con Citibank (un activo de US\$ 14.6 millones), sin una política de cobertura que los enmarque. El precio del crudo tampoco se cubre: solo se amortigua parcialmente con el Fondo de Estabilización de Precios de los Combustibles (FEPC), que alcanza a pocos productos y que en 2025 fue un aporte neto de la empresa (0.43\,\% de los ingresos), no una compensación.
\end{informe}

\section{El objetivo del informe y la hoja de ruta}

El informe sustenta ante la Gerencia General el uso de derivados y responde tres preguntas: (i) ¿qué riesgos financieros son grandes y pueden cubrirse con derivados?; (ii) ¿qué estrategias los cubren y cuánto reducen la pérdida? (un \emph{cross currency swap} con forwards sin entrega para el riesgo cambiario, y un programa de swaps y collares de opciones sobre WTI para el precio del crudo); y (iii) ¿cuánto cuestan, qué riesgo queda y qué condiciones exige su ejecución? Las exposiciones corresponden a los estados financieros auditados al cierre de 2025 y los precios de mercado al 28 de setiembre de 2026. Las cifras están en millones de dólares y las pérdidas, entre paréntesis.

El informe abre con su conclusión, al estilo de una comunicación de un gerente de finanzas a su gerente general: \enquote{PETROPERÚ no tiene colchón para absorber pérdidas de mercado}. Y la cierra con una recomendación \emph{condicional}: aprobar las coberturas solo si los bancos aceptan no pedir garantías en efectivo por encima de cierto monto. Entender por qué esa condición es decisiva es el punto de llegada de este libro.

Cada una de esas expresiones (\emph{cross currency swap}, forward sin entrega, collar, valorizar, garantías) será objeto de un capítulo. Antes necesitamos la herramienta más básica de las finanzas: el valor del dinero en el tiempo.

\begin{cuidado}[confundir riesgo con pérdida]
Decir que \enquote{la empresa tiene riesgo cambiario} no significa que vaya a perder por tipo de cambio. Significa que su resultado depende del tipo de cambio. En 2025, de hecho, PETROPERÚ reportó en el agregado una \emph{ganancia} cambiaria neta de US\$ 13.4 millones, aunque su deuda con el Banco de la Nación aumentó en US\$ 119.3 millones por la apreciación del sol. El riesgo existe aunque el saldo final sea favorable.
\end{cuidado}

\begin{cuidado}[creer que cubrirse siempre gana]
Una cobertura bien diseñada \emph{pierde dinero} en los escenarios en que la posición original gana. Su resultado debe juzgarse junto con la posición que cubre, nunca por separado.
\end{cuidado}

\begin{resumen}
\begin{itemize}
\item PETROPERÚ compra crudo en dólares, vende en soles a precios que siguen al dólar y tiene una deuda importante en soles.
\item Su situación financiera es frágil: pérdidas, capital de trabajo negativo de US\$ 2,296.4 millones, utilidad bruta de US\$ 14.4 millones, calificación en grado especulativo e incertidumbre sobre su continuidad.
\item Su moneda funcional es el dólar, por lo que el sol es una moneda extranjera y sus deudas en soles generan riesgo cambiario.
\item El riesgo financiero es variabilidad del resultado. Se clasifica en riesgo de mercado (tipo de cambio, tasa, crudo), de crédito y de liquidez.
\item Cubrirse crea valor cuando el financiamiento externo es caro: estabiliza la caja para que esté disponible cuando se necesita.
\item Un derivado es un contrato cuyo valor depende de otra variable. Las familias básicas son forwards, futuros, swaps y opciones.
\end{itemize}
\end{resumen}

\section*{Ejercicios}
\addcontentsline{toc}{section}{Ejercicios}

\ejercicio[conceptual] Explique con sus palabras por qué una deuda en soles genera riesgo cambiario para PETROPERÚ, aunque PETROPERÚ sea una empresa peruana.

\ejercicio[cálculo] Con los datos del capítulo, calcule cuántas veces la utilidad bruta de 2025 representa una pérdida de US\$ 33.3 millones en el valor del inventario de crudo. ¿Qué significa ese número?

\ejercicio[cálculo] El sol pasa de 3.40 a 3.70 soles por dólar. Calcule el valor en dólares de una deuda de S/ 3,400 antes y después. ¿Gana o pierde el deudor cuya moneda funcional es el dólar?

\ejercicio[conceptual] Clasifique cada situación como riesgo de mercado, de crédito o de liquidez: (a) un distribuidor mayorista no paga una factura; (b) el WTI cae 20\,\%; (c) la empresa no consigue renovar una línea de crédito que vence; (d) sube la tasa de los bonos del Tesoro de EE.\,UU.

\ejercicio[aplicación] Repita el ejemplo de la caja estable suponiendo que la caja operativa sin cobertura es 140 o 60 con igual probabilidad y que el recargo por financiamiento es 30\,\%. ¿Cuánto ahorra en promedio la cobertura?

\ejercicio[conceptual] ¿Por qué el informe dice que la liquidez \enquote{condiciona el diseño} de las coberturas aunque no se cubra con derivados?
